%% kapa: exam=1 round=25 year=2014 period=2 form=A subject=accounting subject_ko=회계학 seq=7 number=47
%% source: Q-Net 감정평가사 기출문제 (cst00309) artlSeq=5101207 file=extracted/1008/5101207/2014년 제25회 감정평가사 1차시험 기출문제/2교시(A형).hwp
%% conversion: kordoc 4.24.51 -> tools/gen1.py (auto); review=proofread-sonnet; flags=none
\begin{question}{47}
유형자산의 재평가 회계처리에 관한 설명으로 옳은 것은?
\begin{choices}
  \item 재평가는 자산의 장부금액이 공정가치와 중요하게 차이가 나지 않도록 매 보고기간말에 수행한다.
  \item 특정 유형자산을 재평가할 때, 해당 자산이 포함되는 유형자산 분류 전체를 재평가할 필요는 없으며, 개별 유형자산별로 재평가모형을 선택하는 것이 가능하다.
  \item 자산의 장부금액이 재평가로 인하여 증가된 경우에 그 증가액은 동일한 자산에 대하여 이전에 당기손익으로 인식한 재평가감소액이 있다 하더라도 기타포괄손익으로 인식하고 재평가잉여금의 과목으로 자본에 가산한다.
  \item 자산의 장부금액이 재평가로 인하여 감소된 경우에 그 감소액은 당기손익으로 인식한다. 그러나 그 자산에 대한 재평가잉여금의 잔액이 있다면 그 금액을 한도로 재평가감소액을 기타포괄손익으로 인식한다.
  \item 자본에 계상된 재평가잉여금은 그 자산이 제거될 때 이익잉여금으로 대체할 수 있다. 그러나 기업이 그 자산을 사용함에 따라 재평가잉여금의 일부를 대체할 수도 있다.
\end{choices}
\end{question}
\answer{4}
